\label{sec:limitations}
Our approach has important limitations and validity threats.
\paragraph{Construct validity.}
Proxy corpora are noisy mixtures of domains and user populations; alignment should not be interpreted as task success, satisfaction, or usability.
Our overlap metrics measure \emph{distributional similarity of friction themes}, not causal impact.
We partially mitigate this by reporting multiple metrics, sensitivity analyses, and confidence intervals, but the construct remains approximate.

\paragraph{Internal validity.}
LLM simulations can suffer from prompt framing bias and ``algorithmic fidelity'' failures~\cite{amirova2024algorithmicfidelity,gao2024scylla}.
Persona prompts may introduce systematic biases and unrealistic behavior~\cite{li2025personacatch,prpa2024syntheticpersonae}, and prior work shows that
LLM-generated ``synthetic survey'' distributions can diverge from human distributions~\cite{bisbee2024perils}.
We therefore treat simulations as \emph{decision support} and emphasize versioned prompts, schema-constrained outputs, and artifact trails for auditability.

\paragraph{External validity.}
Results depend on the chosen taxonomy, locales, and proxy sources.
While adding OSS issues broadens evidence toward B2B tooling, it remains an indirect proxy and may over-represent power users or issue reporters.
Future work should evaluate domain-specific proxies and compare against controlled studies when available.

\paragraph{Evaluation validity.}
The annotated Amazon subset provides only a weak gold label for category classification and may encode annotator bias.
Even when LLMs are used as evaluators or automatic labelers, they may be biased and inconsistent in rubric application~\cite{liu2023geval,zheng2023mtbench,wang2024notfair}.
We use the labeled subset primarily for calibration and error analysis, not as a definitive performance target.

\paragraph{Grounding heuristic calibration.}
The hallucination proxy ($1-\text{grounding}$) and fabrication proxy rate are computed using a token-overlap heuristic that applies a stop list removing short tokens and common UI terms (e.g., ``button'', ``menu'', ``tab'', ``screen'', ``field'').
Because simulation outputs frequently reference these same UI elements using the same vocabulary, the stop list systematically underestimates true token overlap, yielding hallucination proxy values near 1.0 across all conditions.
These values therefore reflect a conservative lower bound on grounding, not evidence that nearly all output tokens are fabricated; teams should recalibrate the stop list for their UI vocabulary before using this proxy as a go/no-go gate.
The proxy remains useful for \emph{relative} comparisons across agent strategies and prompt versions, where the stop-list bias is consistent.

\paragraph{Cost and model choice.}
Our agent ablation uses Azure-hosted OpenAI models (\texttt{gpt-4.1}, \texttt{gpt-4.1-mini}, \texttt{gpt-5.2}) and the LLM label evaluation uses \texttt{gpt-4.1}; the alignment metrics additionally use a local embedding model (BGE-M3 via Ollama). Stronger paid-model calibration of the alignment pipeline at scale, and a broader sweep over embedding model and subsample size, are deferred to future work to control cost and to avoid iterative overfitting.
